%----------
%	RESEARCH CONTEXT
%----------	

Este Trabajo de Fin de Máster se desarrolla bajo un doble contexto académico e institucional. Por un lado, se desarrolla en contacto con el proyecto IRIS\_IAMEDIA~\cite{irisiamedia}, del que toma el corpus y la metodología, sin formar parte de él. Se trata de un proyecto interdisciplinar en el que colaboran expertas en comunicación e ingenieras e ingenieros expertos en el ámbito de la Inteligencia Artificial~\cite{avi2026}. Aunque el TFM se nutre del ecosistema científico y tecnológico generado por IRIS\_IAMEDIA, ninguno de sus resultados forma parte de los objetivos, entregables ni compromisos oficiales del proyecto, cuya agenda de investigación es independiente. No obstante, la participación de las investigadoras principales (IP) y colaboradoras de IRIS\_IAMEDIA ha permitido disponer de un marco metodológico consolidado, así como de recursos, corpus experimentales y líneas de trabajo en el análisis de sesgos de género que han inspirado y orientado esta propuesta.

Por otro lado, esta investigación se desarrolla tras haber obtenido el primer puesto en la resolución de la VIII Convocatoria Impulsa Visión RTVE de ayudas a la investigación para Trabajos Fin de Máster\footnote{Resolución de la VIII Convocatoria Impulsa Visión RTVE Ayudas a la Investigación para TFM. Disponible en: \url{https://www.rtve.es/rtve/20260304/resultados-viii-impulsa-vision-rtve-ayudas-a-investigacion-tfm-inteligencia-artificial-periodismo-convocatoria/16964519.shtml} [Último acceso: 09/03/2026].}, dentro de la línea de \textit{Inteligencia artificial y periodismo}. Este reconocimiento, otorgado durante la fase de evaluación de la idea y propuesta del proyecto, avala el interés y la necesidad del sector periodístico por contar con este tipo de soluciones. La convergencia de este marco institucional con el respaldo del proyecto IRIS\_IAMEDIA proporciona el escenario idóneo para explorar herramientas tecnológicas aplicables a los desafíos de las redacciones actuales.
